\chapter{Extensions and Open Problems}
\label{ch:open}

The first two sections record consequences of the method that hold without further
work: certification along arbitrary affine lines and polynomial curves, and a comparison
with parallel scalar FRI. The remaining sections state the open problems left by
\cref{ch:corollaries,ch:lower-bounds,ch:parameters}, each in a form that can be attacked
directly.

%=====================================================================
\section{Affine lines and polynomial curves}
%=====================================================================

\begin{definition}[Curve challenges]\label{def:curve-challenges}
Fix $e\ge1$ and public $z=(z_1,\dots,z_k)\in\K^k$. In \cref{prot:sfri} replace the
challenge $r_{i+1}T$ by
\[
\varphi_{i+1}(T):=z_{i+1}+\sum_{t=1}^{e}r_{i+1,t}T^t,\qquad
(r_{i+1,1},\dots,r_{i+1,e})\leftarrow\K^e\ \text{uniform},
\]
raise the degree caps to $d_i^{(e)}:=e\,d_i$, and require $\deg\gt\le ed$. The honest
terminal polynomial is $g(T)=f(\varphi_1(T),\dots,\varphi_k(T))$. The case $e=1$, $z=0$ is
\cref{prot:sfri}; the case $e=1$ certifies $f$ on the affine line
$\{z+cr:c\in\K\}$.
\end{definition}

\begin{proposition}[Curve version]\label{prop:curves}
\Cref{thm:symbolic-list,thm:identification} and all results of \cref{ch:corollaries}
hold for the protocol of \cref{def:curve-challenges}, with $f_G(rT)$ replaced by
$f_G(\varphi(T))$ and $d$ replaced by $ed$. The attacks of \cref{ch:lower-bounds} apply
verbatim, with the same replacements.
\end{proposition}

\begin{proof}
We list the steps of \cref{ch:setting,ch:specialization} that use the form of the
challenge and check each.
\begin{enumerate}[label=\textup{(\roman*)}]
\item \emph{Honest invariant.} \Cref{lem:fold-poly} holds for every $\beta\in\K[T]$, so
folding with $\varphi_{i+1}(T)$ substitutes $Y_{i+1}=\varphi_{i+1}(T)$ in the proof of
\cref{prop:honest}; the $T$-degree of layer $i$ is at most $e\,d_i$.
\item \emph{Specialization.} The proof of \cref{lem:specialization} gives
$\Fold_{\varphi(T)}[w](y)|_{T=c}=\Fold_{\varphi(c)}[w^{(c)}](y)$.
\item \emph{Coupling.} Fix $c\in\K^\times$. Given $\beta_{i+1}$ uniform, the scalar prover
samples $r_{i+1,2},\dots,r_{i+1,e}$ uniformly and sets
$r_{i+1,1}:=c^{-1}\bigl(\beta_{i+1}-z_{i+1}-\sum_{t\ge2}r_{i+1,t}c^t\bigr)$. For fixed
$(c,r_{i+1,2},\dots,r_{i+1,e})$ this is a bijection between $\beta_{i+1}$ and
$r_{i+1,1}$, so the simulated vector $(r_{i+1,t})_t$ is uniform on $\K^e$ and independent
of the history, and $\varphi_{i+1}(c)=\beta_{i+1}$. \Cref{lem:coupling} follows as before.
\item \emph{Root count.} $\gt-f_G(\varphi(T))$ has degree at most $ed$, so
$|Z(\tau)|\le\ell_\alpha ed$.
\item \emph{Injectivity.} For $0\ne h\in\Pa(\K)$, if $h(\varphi(T))=0$ in $\K[T]$ then
$h(\varphi(1))=0$, where $\varphi(1)=z+\sum_tr_{\cdot,t}$ is uniform on $\K^k$; this has
probability at most $\deg h/|\K|\le d/|\K|$ by Schwartz--Zippel. This replaces the
homogeneous-component argument in the proofs of \cref{thm:identification} and
\cref{lem:perturbation}(iv).
\end{enumerate}
No other step depends on the challenge.
\end{proof}

%=====================================================================
\section{Comparison with parallel scalar FRI}
%=====================================================================

The terminal polynomial $g$ has degree $\le d$, so it is also determined by the $d+1$
values $g(c_j)=f(c_jr)$ at public distinct points $c_0,\dots,c_d\in\K^\times$, each of which
is the output of a scalar FRI instance with challenges $c_jr_i$ on the common initial word
$W_0$.

\begin{remark}[Opening cost]\label{rem:parallel}
With leaves grouped by fibre, shared query paths and $W_k$ omitted, the number of opened
$\K$-elements per path (layers $1\le i<k$) is
\[
\text{symbolic: }\ a(k-1)\Bigl(1+\frac{(a-1)k}2\Bigr),\qquad
\text{$d+1$ parallel scalar instances: }\ a(k-1)\bigl(1+(a-1)k\bigr),
\]
so the symbolic protocol opens a fraction
$\frac{1+(a-1)k/2}{1+(a-1)k}\in(\frac12,\frac34]$ of the parallel data; for $a=2$, $k=20$:
$418$ against $798$. In addition, the direct analysis of the parallel protocol bounds
$\Pr[\Acc\wedge\exists j:\gt(c_j)\ne f_\star(c_jr)]$ by a union bound over $j$, which
multiplies the query term by $d+1$; the coupling of \cref{lem:coupling} avoids this loss
because the evaluation point is hidden from the prover.
\end{remark}

%=====================================================================
\section{Mutual correlated agreement beyond the Johnson radius}
%=====================================================================

\begin{problem}\label{prob:mca-beyond}
Let $\rho\in\{\frac12,\frac14,\frac18\}$, $a\ge2$, and
\[
\alpha^\ast\in\Bigl\{\Bigl(\frac{1+\rho}2\Bigr)^2,\ \ \delta_0+\rho\Bigr\}\cap(\rho,\sqrt\rho).
\]
Prove $\eps_{\MCA}(\RS[\K,E,\ell],a,1-\alpha^\ast)\le\mathrm{poly}(n)/|\K|$ and
$\ell_{\alpha^\ast}\le\mathrm{poly}(n)$ for $\K=\F_{p^3}$ and $E=D_{i+1}\subseteq\F_p$.
By \cref{prop:gap-mca}, this closes the gaps for non-equivocation and promise binding.
\end{problem}

\begin{lemma}[Subfield domains are interleavings]\label{lem:subfield-interleave}
Let $E\subseteq\F_p$, $\K=\F_{p^e}$, and fix an $\F_p$-basis $\omega_1,\dots,\omega_e$ of
$\K$. The map
\[
\RS[\F_p,E,\ell]^{e}\to\RS[\K,E,\ell],\qquad (P_1,\dots,P_e)\mapsto\sum_{b=1}^{e}\omega_bP_b,
\]
is an $\F_p$-linear bijection, and a word $w=\sum_b\omega_bw_b:E\to\K$ agrees with
$\sum_b\omega_bP_b$ at $x\in E$ iff $w_b(x)=P_b(x)$ for all $b$.
\end{lemma}

\begin{proof}
Every $P\in\K[X]_{<\ell}$ decomposes uniquely as $\sum_b\omega_bP_b$ with
$P_b\in\F_p[X]_{<\ell}$, and for $x\in\F_p$, $P(x)=\sum_b\omega_bP_b(x)$ with
$P_b(x)\in\F_p$; uniqueness of coordinates in the basis gives the agreement statement.
\end{proof}

By \cref{lem:subfield-interleave}, \cref{prob:mca-beyond} is a statement about
\emph{$e$-fold interleaved} Reed--Solomon codes over the prime field $\F_p$, under the
$\K$-linear generator $\beta\mapsto(1,\beta,\dots,\beta^{a-1})$, whose action mixes the
$e$ coordinates. Two lines of recent work bear on it: deterministic MCA bounds a fixed
number of integer steps beyond the Johnson radius on arbitrary evaluation
sets~\cite{Jo26}, and hidden-derivative methods for list decoding and MCA up to capacity
over prime fields~\cite{DKT26,Jer26}. Results over prime fields apply to the coordinate
codes $\RS[\F_p,E,\ell]$; the missing step is their transfer to the interleaved code with
$\K$-valued challenges. The up-to-capacity versions of the proximity-gap conjectures fail
in general~\cite{BCHKS25}; the thresholds $\alpha^\ast$ above lie strictly above
capacity, so these negative results do not apply directly.

%=====================================================================
\section{Fiat--Shamir}
%=====================================================================

\begin{problem}\label{prob:fs}
Prove round-by-round soundness~\cite{CCHLRRW19} for the binding notions of
\cref{ch:corollaries}, with round errors $\eps_{\MCA}$-type terms for the commit rounds
and the query factor to the power $s$ for the query round, and deduce state-restoration
soundness of the BCS-compiled argument~\cite{BCS16} in the random oracle model.
\end{problem}

The proofs of \cref{ch:path,ch:specialization} have the required shape: every bad event
is attached to a single challenge and depends only on data fixed before that challenge,
and the query round is analysed transcript by transcript. If \cref{prob:fs} is resolved
with these round errors, an adversary making $Q$ oracle queries gains a factor $Q$ on the
query round, and unique binding at $\lambda=128$ with $Q=2^{64}$ requires
\[
s=\Bigl\lceil\frac{\lambda+\log_2Q}{\log_2\frac2{1+\rho}}\Bigr\rceil
=463,\ 284,\ 232,\ 211\qquad(\rho=\tfrac12,\tfrac14,\tfrac18,\tfrac1{16}),
\]
or the corresponding number of proof-of-work bits. Analyses of this type for scalar FRI
are given in~\cite{BGKTTZ23}.

%=====================================================================
\section{Opening cost}
%=====================================================================

\begin{problem}\label{prob:cost}
Construct a variant of \cref{prot:sfri} with the binding guarantees of
\cref{thm:unique-binding} and $O(ak)$ opened $\K$-elements per query path, instead of
$\Theta(a^2k^2)$ (\cref{rem:cost}).
\end{problem}

The quadratic term comes from the $T$-degree $d_i=i(a-1)$ of layer $i$. Checking the
identities at one random point (\cref{rem:single-point}) reduces the verifier's
arithmetic but not the opened data, because the leaves must determine $W_i(x)$ as a
polynomial. Any solution must keep the evaluation point hidden from the prover until the
commitments are fixed (the coupling of \cref{lem:coupling}); revealing it in advance
reduces the protocol to scalar FRI at that point, which binds a single value of $g$.

%=====================================================================
\section{Additive domains}
%=====================================================================

\begin{problem}\label{prob:additive}
Develop the protocol over binary fields, with $D_0$ an affine $\F_2$-subspace and the
maps $\sigma_{i+1}$ given by subspace polynomials of degree $a_{i+1}$, and identify the
terminal polynomial $g(T)$ in terms of the novel polynomial basis~\cite{LCH14,DP24}.
\end{problem}

\Cref{ch:mca,ch:path} use only three properties of the domain tower: fibres of size
$a_{i+1}$, an invertible component map on each fibre (\cref{lem:components}(ii)), and the
identification of folded codewords with Reed--Solomon codewords of the next layer
(\cref{lem:fold-poly}). For additive towers the first two hold with the components
defined by the subspace-polynomial decomposition, and the third holds in the novel
basis. The open point is the analogue of \cref{lem:kronecker} and \cref{prop:honest}: which
polynomial in $T$ the honest fold chain produces, and which multivariate object it
certifies.
